\begin{spacing}{1.5}
\end{spacing}
\vspace{4cm}
\begin{center}
\textbf{{\huge ABSTRACT}}
\end{center}
\vspace{2cm}
\begin{spacing}{1.5}
The rapid growth of research publications on platforms such as arXiv makes it challenging for students and researchers to efficiently extract, understand, and compare key insights from multiple papers. This project presents an AI-based research paper summarization and question answering system implemented using a Retrieval-Augmented Generation (RAG) pipeline. Papers are fetched through the arXiv API, parsed into structured text with section and page-level metadata, chunked, embedded, and indexed in a vector store for semantic retrieval. A Large Language Model generates concise summaries and context-grounded answers to user queries. The system incorporates two key enhancements: responses cite the exact section and page from which information is derived, and it supports side-by-side comparison of two to three papers across aspects such as methodology, datasets, and evaluation. Performance is assessed using a small manually curated question set and qualitative LLM-as-judge criteria, emphasizing relevance, faithfulness, and citation quality.
\end{spacing}
